\documentclass[11pt,a4paper]{article}
\usepackage{curso}
\begin{document}
\portada{Lección}{0.5}{Cadenas de Markov}

\noindent\textbf{Unidad 0.}\quad\textbf{Objetivos:} evolucionar distribuciones con $P$; calcular $\pi$; verificar
balance detallado; relacionar brecha espectral y mezcla. \textbf{Prerrequisitos:} lecciones 0.3 y 0.4.

\section{¿Por qué importa?}
Gibbs, Metropolis, un arreglo de pbits y una partícula que salta entre pozos son cadenas de Markov: el próximo
estado depende solo del actual. Toda la Unidad 4 (muestreo), el hardware de la Unidad 5 y la búsqueda de orden
por difusión de la Unidad 7 descansan en esta lección \cite{levin}.

\section{Definiciones}
\begin{definicion}[ (cadena de Markov finita)]
Una sucesión $x_0,x_1,\dots$ en $\{1,\dots,K\}$ con $P(x_{n+1}=j\mid x_n=i,x_{n-1},\dots)=P_{ij}$. $P$ es
\emph{estocástica por filas}: $P_{ij}\ge0$, $\sum_jP_{ij}=1$.
\end{definicion}
Si $p_n$ (vector \emph{fila}) es la distribución de $x_n$, entonces $p_{n+1}=p_nP$ y $p_n=p_0P^n$.

\section{Distribución estacionaria}
$\pi$ es estacionaria si $\pi P=\pi$, $\pi_i\ge0$, $\sum\pi_i=1$: vector propio izquierdo con valor propio 1.
\begin{teorema}[ergódico para cadenas finitas]
Si $P$ es irreducible (todo estado alcanza a todo estado) y aperiódica, existe una única $\pi$ y
$p_n\to\pi$ para todo $p_0$. Además, las fracciones de tiempo de visita convergen a $\pi$ casi seguramente.
\end{teorema}

\section{Balance detallado y reversibilidad}
\begin{definicion}[ (balance detallado)]
$\pi_iP_{ij}=\pi_jP_{ji}$ para todo $i,j$.
\end{definicion}
\begin{proposicion}
Balance detallado $\Rightarrow$ estacionariedad: $\sum_i\pi_iP_{ij}=\sum_i\pi_jP_{ji}=\pi_j\sum_iP_{ji}=\pi_j$.
\end{proposicion}
El recíproco es falso: la cadena cíclica de 3 estados con probabilidad 0,8 hacia adelante, 0,1 hacia atrás y 0,1
de quedarse tiene $\pi$ uniforme pero una corriente neta $\frac13(0{,}8-0{,}1)$ en cada arista. Las cadenas con
balance detallado (reversibles) describen sistemas en \emph{equilibrio termodinámico}; las otras, estados
estacionarios fuera del equilibrio (Unidad 3). \textbf{Receta:} dados pesos simétricos $W_{ij}=W_{ji}\ge0$,
$P_{ij}=W_{ij}/\sum_kW_{ik}$ es reversible con $\pi_i\propto\sum_kW_{ik}$ (así se construyó el ejemplo).

\section{Brecha espectral y tiempo de mezcla}
Para $P$ reversible, $D^{1/2}PD^{-1/2}$ con $D=\diag(\pi)$ es simétrica: los valores propios son reales,
$1=\lambda_1>\lambda_2\ge\dots\ge\lambda_K>-1$. Sea $\lambda_*=\max(|\lambda_2|,|\lambda_K|)$. Entonces
\begin{equation}\|p_n-\pi\|_1\le C\,\lambda_*^n,\qquad t_{\text{mix}}(\varepsilon)\lesssim\frac{1}{1-\lambda_*}\ln\frac{1}{\varepsilon\,\pi_{\min}}.\end{equation}
$1-\lambda_*$ es la \emph{brecha espectral} (absoluta). En modelos de Ising cerca de la temperatura crítica, la brecha
de Gibbs se cierra y el número de barridos necesario se dispara (Unidad 4.4).

\begin{ejemplo}
$P=\begin{pmatrix}2/3&1/3&0\\1/4&1/2&1/4\\0&1/3&2/3\end{pmatrix}$, construida con $W=\begin{pmatrix}2&1&0\\1&2&1\\0&1&2\end{pmatrix}$.
Entonces $\pi\propto(3,4,3)$, es decir $\pi=(0{,}3;0{,}4;0{,}3)$. Balance detallado: $\pi_1P_{12}=0{,}3\cdot\frac13=0{,}1=0{,}4\cdot\frac14=\pi_2P_{21}$
(y por simetría entre 2 y 3). Valores propios: $1,\frac23,\frac16$; brecha $\frac13$; la distancia a $\pi$ se reduce
asintóticamente en un factor $\frac23$ por paso, como muestra la recta en escala logarítmica.
\end{ejemplo}

\begin{aplicacion}
Gibbs y Metropolis (Unidad 4) se diseñan para que $\pi$ sea la distribución de Boltzmann y cumplen balance
detallado. Un TSU es una cadena de Markov física: su brecha espectral decide los barridos necesarios. El
laboratorio L1 pide: cadenas de 3 y 50 estados; $\pi$ por iteración de potencias y por descomposición espectral;
gráfica de $\|p_n-\pi\|_1$ frente a $n$ comparada con $\lambda_2^n$.
\end{aplicacion}

\section{Conexiones}
Lección 0.6 (Fourier en $\Z/r\Z$ diagonaliza caminatas en ciclos), 1.2 (Boltzmann), 2.3 (Fokker--Planck, versión
continua), 2.8 (tiempo continuo y Gillespie), 4.3--4.4 (MCMC y diagnósticos), 7.3 (caminatas en grafos de Cayley).

\begin{error}
\begin{itemize}[nosep]
\item Filas frente a columnas: con $p$ fila, $p_{n+1}=p_nP$ y las filas de $P$ suman 1.
\item Estacionaria $\not\Rightarrow$ balance detallado.
\item Cadena periódica: tiene $\pi$ pero $p_n$ no converge (oscila).
\item La mezcla la decide $\lambda_*=\max|\lambda_{i\ge2}|$, no solo $\lambda_2$.
\end{itemize}
\end{error}

\section{Resumen}
$p_{n+1}=p_nP$; $\pi P=\pi$; balance detallado $\Rightarrow$ estacionariedad y equilibrio; $\|p_n-\pi\|\lesssim C\lambda_*^n$.

\begin{thebibliography}{9}\small
\bibitem[3]{levin} D. A. Levin, Y. Peres y E. L. Wilmer, \emph{Markov Chains and Mixing Times}, AMS, 2009.
\bibitem[13]{katz} H. G. Katzgraber, Introduction to Monte Carlo methods, 2009. \url{https://arxiv.org/abs/0905.1629}
\end{thebibliography}
\end{document}
